\documentclass[11pt]{article}

\usepackage[utf8]{inputenc}
\usepackage[T1]{fontenc}
\usepackage[margin=1in]{geometry}
\usepackage{graphicx}
\usepackage{amsmath,amssymb}
\usepackage{booktabs}
\usepackage{array}
\usepackage{enumitem}
\usepackage{tcolorbox}
\usepackage{float}
\usepackage{placeins}
\usepackage[hidelinks]{hyperref}
\usepackage{textcomp}
\usepackage{longtable}
\usepackage{xcolor}

\newcommand{\Analog}{\textsc{Analog}}
\newcommand{\None}{\textsc{None (Analog)}}
\newcommand{\Alternating}{\textsc{Alternating}}
\newcommand{\FullSpiking}{\textsc{Full Spiking}}
\newcommand{\HybridLate}{\textsc{Hybrid Late}}
\newcommand{\HybridEarly}{\textsc{Hybrid Early}}
\newcommand{\HybridMid}{\textsc{Hybrid Mid}}

\newtcolorbox{keypoint}{%
  colback=black!3, colframe=black!60, boxrule=0.6pt, arc=1pt,
  left=8pt, right=8pt, top=6pt, bottom=6pt
}
\newtcolorbox{supersede}{%
  colback=yellow!8, colframe=orange!70!black, boxrule=0.8pt, arc=1pt,
  left=8pt, right=8pt, top=6pt, bottom=6pt
}

\title{\Large Rotation Invariance in Spiking and Analog\\
$p4$-Equivariant CNNs\\[4pt]
\large Consolidated Findings Report}
\author{Sahil Singh\\
CMI Summer Research Internship\\
Supervisor: Prof. K.V. Subrahmanyam}
\date{August 2026}

\begin{document}

\maketitle

\begin{abstract}
This report consolidates four previously separate reports and two post-hoc corrections into one
internally consistent record, spanning a full-capacity (7-layer, 10-channel, 24{,}744-parameter) $p4$-equivariant
spiking/analog CNN and a reduced-capacity (5-layer, 4-channel, 2{,}420-parameter) companion model run
under two independent training regimes, plus a follow-up ``Priority Experimental Programme'' (Experiments
2--4: spike placement, layer/time representation distance, and Fourier analysis of the rotational response).
No new experiments are run and no model is retrained here --- every table and figure reads from results
already produced across the four source reports, plus two corrections applied to the analysis code
after a code-vs-methodology audit uncovered two genuine metric-design issues (a fixed epsilon
inappropriate for a narrow layer, and single-draw Monte-Carlo noise in a stochastic spike-encoding
comparison). The findings that survive scrutiny at every capacity and regime tested: a
depth-wise representational collapse at the deepest layer(s), present even in the zero-spiking control,
identifying it as an architectural rather than spiking signature; a rise-then-plateau (not decay)
temporal rotation-sensitivity pattern, contradicting this project's original temporal-averaging
hypothesis; a clean, reproducible relationship between spiking depth and how many timesteps are needed
before rotation-robust accuracy appears; and, in Experiment 4, the cleanest confirmation in the project
of the $p4$ architecture's exact discrete-symmetry guarantee. Two findings required real correction on
top of that: the reduced-capacity model's group-angle equivariance error, and its layer/time
representation-distance metric at the narrowest layer, both of which were substantially distorted by
metric-design artifacts rather than reflecting the underlying model. Both are corrected here and the
corrected numbers explicitly supersede the originals.
\end{abstract}

\tableofcontents

% ===========================================================================
\section{How to Read This Report}
\label{sec:part0}

This document consolidates four previously separate reports and two post-hoc corrections into a
single, internally consistent record. It reproduces no new experiments and trains no new models ---
every table and figure below reads from results already on disk in \texttt{\textasciitilde/snn/} at the
time of writing.

\paragraph{The four source reports being merged.}
\begin{itemize}[leftmargin=1.6em]
\item \textbf{Report 1} (``full-capacity companion report''): 7-layer, 10-channel, 24{,}744-parameter
\texttt{SpikingP4CNNSmall}, $T=50$ simulation timesteps, 20 training epochs, 4 spike-placement
configurations, 3 seeds. File: \texttt{findings\_report.tex} / \texttt{findings\_report.pdf}.
\item \textbf{Report 2} (``reduced-capacity companion report''): 5-layer, 4-channel, 2{,}420-parameter
\texttt{SpikingP4CNNSmall}, run under two independent regimes ($T{=}50$/10ep and $T{=}30$/5ep), 5
spike-placement configurations, 3 seeds. Files: \texttt{findings\_report\_small.tex} (both regimes) and
\texttt{findings\_report\_t30.tex} ($T{=}30$ regime, revised after seed 1 arrived, more cautious framing
of the group-angle-error finding).
\item \textbf{Report 3} (``Priority Experimental Programme, Experiments 2--4''): Experiment 2 (spike
placement vs.\ accuracy/$E_{90}$), Experiment 3 (layer/time representation distance $D_{\ell,t}$), and
Experiment 4 (Fourier decomposition of the rotation response), run on the same reduced-capacity
checkpoints as Report 2, both regimes, no retraining. File: \texttt{findings\_report\_experiments234.tex}.
\item \textbf{The corrections.} A code-vs-methodology audit of three specific anomalies raised across
Reports 2 and 3, followed by two concrete fixes applied to the analysis code: Fix 1 (a layer-relative
epsilon and floor-masking for the $D_{\ell,t}$ metric) and Fix 2 (multi-draw Monte-Carlo averaging for
the $E_{90}$ direct-equivariance-error statistic). No model, training, or architecture code was touched
by either fix --- both are corrections to analysis/metric code only.
\end{itemize}

\begin{supersede}
\textbf{Superseding notice.} The $E_{90}$ numbers in Section~\ref{sec:part5} and the $D_{\ell,t}$
numbers in Section~\ref{sec:part6} \emph{supersede} the equivalent single-draw / fixed-epsilon numbers
reported in Reports 2 and 3. Every place this applies is flagged inline with the word ``supersedes''
rather than left implicit. The mechanism behind both corrections is explained once, in
Section~\ref{sec:part1}'s closing subsection, and referenced from Sections~\ref{sec:part5}/\ref{sec:part6}
as ``the correction described in Section~\ref{sec:part1}'' rather than re-explained.
\end{supersede}

\paragraph{Reading order.} Sections~\ref{sec:part1}--\ref{sec:part2} establish shared setup and the
accuracy backdrop against which every other finding should be read. Sections~\ref{sec:part3}--\ref{sec:part8}
present the substantive findings, one experiment family per section, each explicitly stating which
report(s) and which regime(s) it draws from. Section~\ref{sec:part9} is a permanent settled-questions
record of the code-vs-methodology audit, so this project's anomalies are not re-investigated by a future
reader. Section~\ref{sec:part10} synthesizes at the confidence level each finding actually supports.
Section~\ref{sec:part11} consolidates limitations and next steps against the current deadlines.

\FloatBarrier

% ===========================================================================
\section{Setup}
\label{sec:part1}

\subsection{Architecture}

Both variants share the same qualitative design: a stack of $p4$-equivariant convolutions
(\texttt{P4ConvZ2} lifting the input into the $p4$ group, then \texttt{P4ConvP4} layers operating within
it), \texttt{P4BatchNorm2d} (sharing statistics and affine parameters across the 4 orientation channels
--- see the callout below), a final group-pooling step (max over the 4 orientation channels) producing
an orientation-invariant feature map, and a linear classifier. Full implementation:
\texttt{gconv\_small.py} (\texttt{P4ConvZ2} / \texttt{P4ConvP4} / \texttt{P4BatchNorm2d} /
\texttt{P4GroupPool}) and \texttt{p4cnn\_small.py} / \texttt{spiking\_p4cnn\_small.py} (the analog and
spiking network definitions).

\begin{table}[H]
\centering
\begin{tabular}{@{}lll@{}}
\toprule
 & \textbf{Full-capacity variant} & \textbf{Reduced-capacity variant} \\
\midrule
Layers & 7 & 5 \\
Channels per layer & 10 & 4 \\
Parameters & 24{,}744 & 2{,}420 ($10.2\times$ fewer) \\
Kernels & six $3\times3$ + one $4\times4$ (valid) & $3\times3$ throughout, two $2\times2$ maxpools \\
Used in & Report 1 only & Reports 2 and 3 \\
\bottomrule
\end{tabular}
\caption{Architecture comparison.}
\end{table}

\begin{keypoint}
\texttt{P4BatchNorm2d} exists because standard BatchNorm assigns each of the 4 orientation channels
independent running statistics, which breaks the architecture's equivariance guarantee once trained
upright-only (verified: untrained equivariance error $\sim 10^{-6}$; after training with standard
BatchNorm, $\sim 0.1$--$0.2$). Sharing statistics and affine parameters across the 4 orientation channels
restores equivariance to $\sim 10^{-6}$ even after training. This fix is used throughout every result in
this document --- it is a settled, resolved issue, listed here for completeness and cross-referenced
from Section~\ref{sec:part9}.
\end{keypoint}

\subsection{Task}
4-class rotated-MNIST: digits $\{3,4,8,9\}$, chosen because none is rotationally self-symmetric under
$90^\circ/180^\circ$. Models are trained on upright ($0^\circ$) images only and evaluated on a 24-angle
sweep, $\theta \in \{0^\circ, 15^\circ, \ldots, 345^\circ\}$, on the held-out test split (Reports 1--2)
or a fixed probe set (Report 3 --- see Sections~\ref{sec:part5}--\ref{sec:part7} setup notes).

\subsection{Spike-Placement Configurations and Which Report Trained Which}

\begin{table}[H]
\centering
\small
\begin{tabular}{@{}llll@{}}
\toprule
\textbf{Configuration} & \textbf{Spiking layers} & \textbf{Full-cap.\ (Report 1)} & \textbf{Reduced-cap.\ (Reports 2--3)} \\
\midrule
Analog / None & none & yes (``Analog'') & yes (``None'') \\
Alternating & interleaved (2,4,6 full-cap; 2,4 reduced) & yes & yes \\
Full Spiking & all layers & yes & yes \\
Hybrid Late & final layer only & yes & yes \\
Hybrid Early & first 1--2 layers & not trained --- see note & yes \\
Hybrid Mid & middle layers & not trained --- see note & does not exist at depth 5 \\
\bottomrule
\end{tabular}
\caption{Configuration coverage by report.}
\end{table}

\noindent\textit{Note: Report 1 explicitly states \HybridEarly\ and \HybridMid\ were implemented in the
codebase but not yet trained under its corrected setup, so Report 1 says nothing about them. Both were
trained for the reduced-capacity model (Reports 2--3); \HybridMid\ has no reduced-capacity analogue
since a 5-layer network has no distinct ``middle'' block the way a 7-layer network does.}

\subsection{Training Regimes and Seed Counts}

\begin{table}[H]
\centering
\small
\begin{tabular}{@{}llccccp{4.1cm}l@{}}
\toprule
\textbf{Regime} & \textbf{Model} & $T$ & \textbf{Epochs} & \textbf{Params} & \textbf{Seeds} & \textbf{Repr.\ analysis coverage} & \textbf{Used by} \\
\midrule
Full-capacity & 7L/10ch & 50 & 20 & 24{,}744 & 3 & all 3 seeds & Report 1 \\
Regime A & 5L/4ch & 50 & 10 & 2{,}420 & 3 & all 3 seeds & Reports 2, 3 \\
Regime B & 5L/4ch & 30 & 5 & 2{,}420 & 3 & seeds 0--1 only$^{*}$ & Reports 2, 3 \\
\bottomrule
\end{tabular}
\caption{Training regimes. All regimes: Poisson/Bernoulli rate coding, Adam, lr $10^{-3}$, batch 64,
best checkpoint by $0^\circ$ test accuracy, \texttt{gpuserver.cmi.ac.in}. $^{*}$Seed 2's
representation-level analysis was intentionally left incomplete, by explicit instruction (all 3 seeds
were trained to completion in every regime).}
\end{table}

\subsection{The Two Corrections (Referenced Throughout as ``the Correction Described in Section~\ref{sec:part1}'')}
\label{sec:corrections}

Both corrections were produced by a code-vs-methodology audit that investigated three specific
anomalies raised in Reports 2--3 (full detail and verdicts in Section~\ref{sec:part9}). Two of the three
anomalies turned out to be genuine methodological artifacts worth correcting in the analysis code ---
not bugs in the model, training, or the $p4$-equivariant conv/BatchNorm implementation, which was
independently verified to have no mode-dependent branching.

\paragraph{Fix 1 ($D_{\ell,t}$ epsilon).} A single fixed epsilon ($10^{-8}$) is added to a reference-norm
denominator when computing the layer/time representation distance $D_{\ell,t}(\theta)$ (Report 3,
Experiment 3). At the network's narrowest layer (layer 5, 16 units), the true reference norm is
legitimately exactly zero for a nontrivial fraction of (timestep, image) cells --- dividing by the
$10^{-8}$ floor rather than a real signal produced the reported $\sim 10^7$ blowups. Fix: a per-layer
epsilon ($\max(10^{-8}, 1\%$ of the median nonzero reference norm at that layer$)$, computed once from
the reference batch) plus explicit exclusion of any cell whose raw reference norm falls below $10^{-4}$
from the reported mean. Applied in \texttt{representations\_small.py}'s \texttt{layer\_time\_distance()}
and rerun via the new \texttt{plot\_experiment3\_layertime\_fixed.py}.

\paragraph{Fix 2 ($E_{90}$ multi-draw averaging).} The direct equivariance-error computation ($E_{90}$,
used for Experiment 2's ``does spike placement matter'' ordering claim) resets the encoding RNG to the
same seed value for both the upright image $x$ and its rotation $R_\theta x$ before drawing the
Poisson/Bernoulli spike encoding. Because rotation moves image content between tensor positions, sharing
a seed value reproduces the same raw random stream but does not reproduce a matched spike encoding
under rotation --- it only gives run-to-run reproducibility, not the variance cancellation the
convention seems to have intended. A single draw's $E_{90}$ is therefore dominated by Monte-Carlo
spike-sampling variance, not a stable property of the trained model (confirmed empirically: re-drawing
the encoding seed alone, with no rotation or config change, reproduces the reported group-angle
magnitudes and ordering). Fix: draw the encoding 8 independent times per (model seed, config), each with
a genuinely different encoding seed ($\text{seed}\times1000 + \text{draw\_index}$, never reused), and
report the mean and std across draws instead of a single value. Added as
\texttt{equivariance\_error\_multidraw()} in \texttt{representations\_small.py} and run via the new
\texttt{plot\_experiment2\_summary\_multidraw.py}.

\textit{Both fixes touch analysis/metric code only (\texttt{representations\_small.py} plus two new
driver scripts); \texttt{gconv\_small.py}, \texttt{p4cnn\_small.py}, \texttt{spiking\_p4cnn\_small.py},
and every training script are untouched, and no model was retrained.}

\FloatBarrier

% ===========================================================================
\section{Baseline Accuracy (All Configs, All Regimes)}
\label{sec:part2}

Across every regime in this project, the qualitative story is the same: once training is done properly
(\texttt{P4BatchNorm2d}, adequate epochs), rotation-accuracy curves for every spike-placement
configuration sit within a few percentage points of each other, with substantial seed-to-seed error-bar
overlap at the hardest angles. Accuracy alone does not distinguish spike placement in any regime tested.

\subsection{Full Capacity (Report 1): $T{=}50$, 20 Epochs, 24{,}744 Params}

\textit{Report 1 does not publish a per-configuration numeric accuracy table --- only the comparison
figure below plus qualitative text are available in that report. Its text states the four
configurations' curves sit within a few points of each other at every angle, with substantial
error-bar overlap at the hardest angles ($45^\circ,135^\circ,225^\circ,315^\circ$), confirming that two
dramatic accuracy-level findings from an earlier, compute-constrained ($T{=}30$, 5-epoch, single-seed,
CPU-only) draft of the same project did not replicate under this corrected setup (\FullSpiking's earlier
flat 51--57\% and \HybridLate's earlier collapse to chance at non-group angles). Under the corrected
setup, \FullSpiking\ reaches $\sim 99\%$ upright and $\sim 94\%$ mean accuracy across all angles;
\HybridLate's worst case across all seeds and angles is 80.9\%.}

\begin{figure}[H]
\centering
\includegraphics[width=0.85\textwidth]{plots/n4_comparison_final.png}
\caption{Full capacity: 4-class rotation-angle sweep, mean $\pm$ std over 3 seeds,
\Analog\ / \Alternating\ / \FullSpiking\ / \HybridLate\ (\HybridEarly/\HybridMid: not trained, see
Section~1.3). \texttt{plots/n4\_comparison\_final.png}}
\end{figure}

\FloatBarrier

\subsection{Reduced Capacity, Regime A: $T{=}50$, 10 Epochs, 2{,}420 Params}

\begin{table}[H]
\centering
\begin{tabular}{@{}lrr@{}}
\toprule
\textbf{Configuration} & Best $0^\circ$ acc.\ (\%), mean$\pm$std, 3 seeds & Mean acc.\ over 24 angles (\%) \\
\midrule
\None            & $97.67 \pm 0.48$ & $84.77$ \\
\Alternating     & $97.39 \pm 0.73$ & $85.62$ \\
\FullSpiking     & $96.39 \pm 0.36$ & $86.30$ \\
\HybridEarly     & $95.14 \pm 3.53$ & $81.32$ \\
\HybridLate      & $97.60 \pm 0.24$ & $85.84$ \\
\bottomrule
\end{tabular}
\caption{Regime A baseline accuracy.}
\end{table}

\begin{figure}[H]
\centering
\includegraphics[width=0.85\textwidth]{plots_small/small_n5_comparison.png}
\caption{Regime A ($T{=}50$, 10 epochs), all 5 configurations, mean $\pm$ std over 3 seeds.
\texttt{plots\_small/small\_n5\_comparison.png}}
\end{figure}

\FloatBarrier

\subsection{Reduced Capacity, Regime B: $T{=}30$, 5 Epochs, 2{,}420 Params}

\begin{table}[H]
\centering
\begin{tabular}{@{}lrr@{}}
\toprule
\textbf{Configuration} & Best $0^\circ$ acc.\ (\%), mean$\pm$std, 3 seeds & Mean acc.\ over 24 angles (\%) \\
\midrule
\None            & $97.25 \pm 0.34$ & $84.52$ \\
\Alternating     & $96.71 \pm 0.12$ & $86.87$ \\
\FullSpiking     & $95.14 \pm 0.53$ & $84.31$ \\
\HybridEarly     & $91.39 \pm 5.71$ & $77.86$ \\
\HybridLate      & $97.09 \pm 0.08$ & $82.03$ \\
\bottomrule
\end{tabular}
\caption{Regime B baseline accuracy.}
\end{table}

\begin{figure}[H]
\centering
\includegraphics[width=0.85\textwidth]{plots_small_t30/small_n5_comparison.png}
\caption{Regime B ($T{=}30$, 5 epochs), all 5 configurations, mean $\pm$ std over 3 seeds.
\texttt{plots\_small\_t30/small\_n5\_comparison.png}}
\end{figure}

\FloatBarrier

One configuration breaks the ``accuracy is uninformative'' pattern in both reduced-capacity regimes:
\HybridEarly\ shows a conspicuously large seed-to-seed standard deviation ($\pm 3.53$ at Regime A,
$\pm 5.71$ at Regime B), driven by a single unstable seed (seed 2: 90.20\% and 83.33\%
best-upright-accuracy respectively, versus $\geq 95\%$ for seeds 0 and 1 in both regimes). This
instability has no full-capacity analogue since Report 1 never trained \HybridEarly; it is plausibly a
capacity effect (at only 4 channels, \HybridEarly's spiking layers 1--2 sit immediately after the
lifting convolution, where the network has the least redundancy to absorb a bad initialization's
interaction with the surrogate-gradient spiking nonlinearity) and was not investigated further
(Section~\ref{sec:part11}).

\FloatBarrier

% ===========================================================================
\section{Layer-wise Representation Similarity (Experiment 5.1)}
\label{sec:part3}

This finding was consistent everywhere it was tested --- full capacity, both reduced-capacity regimes,
every spike-placement configuration including the zero-spiking control --- and is stated here plainly,
without hedging, because the evidence does not call for any.

\begin{keypoint}
Representational similarity to the upright image (measured by cosine similarity, CKA, and SVCCA between
$h_\ell(x)$ and $h_\ell(R_\theta x)$) stays high through the early and middle layers and collapses only
at the final layer(s), in every configuration tested at every capacity, including \Analog/\None, which
has zero spiking neurons anywhere in the network. Whatever causes deep-layer representations to become
rotation-sensitive is a property of the $p4$-equivariant architecture itself (most plausibly the
group-pooling and final linear-classifier stages, the only components not built from strictly
equivariant convolutions) --- not something introduced by spiking dynamics.
\end{keypoint}

\subsection{Full Capacity (Report 1)}

CKA similarity to the upright representation stays at or above roughly 0.97--0.99 through layers 1--5
in every configuration, then drops sharply in layers 6--7 to roughly 0.70--0.85 at the hardest angles
(typically $15^\circ,30^\circ,45^\circ,60^\circ$ etc., the non-group angles nearest a group angle).
Spike placement modulates the shape somewhat --- \Alternating\ and \FullSpiking\ show a slightly earlier
onset of divergence (visible already around layer 5) relative to \Analog\ and \HybridLate\ --- but not
the existence or final location of the collapse: all four configurations land in the same 0.70--0.85
CKA range by layer 7 at the hardest angles. Cosine similarity is much lower and noisier throughout
(typically 0.3--0.6 even at shallow layers) since it demands exact neuron-by-neuron agreement with no
tolerance for internal reshuffling of which channel encodes which feature; CKA and SVCCA both stay near
1.0 through layer 5 and diverge only at layers 6--7.

\begin{figure}[H]
\centering
\includegraphics[width=0.75\textwidth]{plots/analog_layerwise_cka.png}
\caption{Full capacity, \Analog\ --- CKA similarity to upright vs.\ layer, one line per test angle.
\texttt{plots/analog\_layerwise\_cka.png}}
\end{figure}
\begin{figure}[H]
\centering
\includegraphics[width=0.75\textwidth]{plots/alternating_layerwise_cka.png}
\caption{Full capacity, \Alternating\ --- same quantity. \texttt{plots/alternating\_layerwise\_cka.png}}
\end{figure}
\begin{figure}[H]
\centering
\includegraphics[width=0.75\textwidth]{plots/full_layerwise_cka.png}
\caption{Full capacity, \FullSpiking\ --- same quantity. \texttt{plots/full\_layerwise\_cka.png}}
\end{figure}

\FloatBarrier
\textit{\HybridLate's standalone aggregate CKA plot (\texttt{hybrid\_late\_layerwise\_cka.png}) was not
generated in Report 1's pass --- explicitly noted there, not silently omitted here. Its per-seed 5.1
figures were used instead in that report and show the same collapse pattern.}

\subsection{Reduced Capacity, Both Regimes (Reports 2--3)}

With only 5 layers instead of 7 there is less ``room'' for a late-layer-only collapse to look visually
distinct from a whole-network effect, yet the qualitative pattern reproduces cleanly: similarity to the
upright representation stays high through the early-to-middle layers and drops at the final layer(s),
in every configuration and both training regimes, including \None. This is the single most important
replication in Report 2: the architectural (not spiking) origin of the collapse holds at $10.2\times$
less capacity, 40\% less simulation time (Regime B), and a quarter of the training epochs (Regime B).

\begin{figure}[H]
\centering
\includegraphics[width=0.75\textwidth]{plots_small/repr_seed0/similarity_comparison_cka.png}
\caption{Regime A, seed 0 --- CKA similarity to upright vs.\ layer, all 5 configurations overlaid.
\texttt{plots\_small/repr\_seed0/similarity\_comparison\_cka.png}}
\end{figure}
\begin{figure}[H]
\centering
\includegraphics[width=0.75\textwidth]{plots_small_t30/repr_seed0/similarity_comparison_cka.png}
\caption{Regime B, seed 0 --- same comparison.
\texttt{plots\_small\_t30/repr\_seed0/similarity\_comparison\_cka.png}}
\end{figure}

\FloatBarrier
\textit{Spike placement's secondary modulation of exactly where the divergence first becomes visible
(\Alternating/\FullSpiking\ diverging one layer earlier than \None/\HybridLate, per the full-capacity
finding) is harder to isolate cleanly with only 5 layers to distribute the effect across --- this is a
genuine capacity-driven limitation on how finely the secondary claim can be tested at this depth, not a
contradiction of it.}

% ===========================================================================
\section{Rotation Sensitivity, by Layer and by Time (Experiments 5.3, 5.4)}
\label{sec:part4}

\subsection{5.3: Distribution of $\mathcal{R}_i$ Across Layers}

$\mathcal{R}_i = \mathrm{Var}_\theta(a_i(R_\theta x))$, the rotation sensitivity of neuron $i$ over the
24 test angles. In every configuration and every capacity, the final layer sits at or near the top of
the $\mathcal{R}_i$ ranking, consistent with Section~\ref{sec:part3}'s identification of the last
layer(s) as where invariance breaks down. The shape of the climb to the final layer differs
meaningfully by spike placement, at full capacity:

\begin{table}[H]
\centering
\begin{tabular}{@{}lp{10.5cm}@{}}
\toprule
\textbf{Configuration} & \textbf{Shape of $\mathcal{R}_i$ climb to final layer (full capacity)} \\
\midrule
\Alternating & sharp, order-of-magnitude jump only at the final layer: median $\mathcal{R}_i \sim 10$--100
for the earlier layers, jumping to $\sim 400$--800 at layer 7 \\[3pt]
\FullSpiking & gradual, full-depth climb: median $\mathcal{R}_i$ rises roughly monotonically from
$\sim 10$ (layer 1) to $\sim 90$--100 (layer 7), no single sharp jump \\[3pt]
\HybridLate & nearly flat across all layers: median $\mathcal{R}_i$ in a comparatively narrow band
($\sim 10$--80) at every depth, including the final layer \\[3pt]
\Analog & gradual full-depth climb similar in shape to \FullSpiking, but $\sim 1000\times$ smaller in
absolute scale ($\sim 0.01$--0.05 throughout) \\
\bottomrule
\end{tabular}
\caption{Per-configuration $\mathcal{R}_i$ shape, full capacity.}
\end{table}

\begin{figure}[H]
\centering
\includegraphics[width=0.75\textwidth]{plots/full_rotation_sensitivity.png}
\caption{Full capacity, \FullSpiking\ --- $\mathcal{R}_i$ distribution by layer (log scale).
\texttt{plots/full\_rotation\_sensitivity.png}}
\end{figure}

\FloatBarrier
\textit{\HybridLate's standalone aggregate boxplot (\texttt{hybrid\_late\_rotation\_sensitivity.png})
was likewise not generated in Report 1's pass; its per-seed 5.3 figures were used instead there.}

\textit{Reduced capacity: the same qualitative fingerprint (sharp late jump for \Alternating, gradual
climb for \FullSpiking/\None, flatter profile for \HybridLate) reproduces, but is harder to resolve
cleanly with only 5 layers to distribute the pattern across. This is consistent with, not contradicted
by, the full-capacity finding --- the reduced-capacity report is underpowered to cleanly distinguish the
shapes at this depth, not evidence the shapes actually converge.}

\begin{figure}[H]
\centering
\includegraphics[width=0.75\textwidth]{plots_small/repr_seed0/rotsens_median_comparison.png}
\caption{Regime A, seed 0 --- median $\mathcal{R}_i$ per layer, all 5 configurations overlaid (log scale).
\texttt{plots\_small/repr\_seed0/rotsens\_median\_comparison.png}}
\end{figure}

\FloatBarrier
\begin{keypoint}
$\mathcal{R}_i$'s absolute scale depends on whether the underlying activation is a continuous ReLU
output or a spike count, so the fair cross-configuration comparison is the \emph{shape} of the
$\mathcal{R}_i$-vs-layer curve, not the raw magnitude. Within a single configuration, comparisons across
layers or seeds are on safe footing.
\end{keypoint}

\subsection{5.4: Rotation Sensitivity Through Time}

$\mathcal{R}_i(t)$: the same quantity computed from the instantaneous per-timestep spike output alone,
plotted against timestep $t$. In every configuration, every seed, and every capacity/regime tested,
$\mathcal{R}_i(t)$ rises quickly over the first few timesteps and then plateaus for the remainder of the
simulation, with no visible downward trend. This directly contradicts the temporal-contraction
hypothesis that motivated this experiment (that temporal averaging should progressively reduce
sensitivity, $\mathcal{R}_i(t)$ decaying toward 0 as spikes accumulate) --- sensitivity is set almost
immediately and then stays fixed.

\begin{figure}[H]
\centering
\includegraphics[width=0.75\textwidth]{plots/full_seed0/full_seed0_5.4_rotation_sensitivity_through_time.png}
\caption{Full capacity, \FullSpiking, seed 0 --- $\mathcal{R}_i(t)$ vs.\ timestep, all 7 layers.
\texttt{plots/full\_seed0/full\_seed0\_5.4\_rotation\_sensitivity\_through\_time.png}}
\end{figure}
\begin{figure}[H]
\centering
\includegraphics[width=0.75\textwidth]{plots_small/repr_seed0/rotsens_temporal_full.png}
\caption{Regime A, \FullSpiking, seed 0 --- median $\mathcal{R}_i(t)$ vs.\ timestep, all 5 layers.
\texttt{plots\_small/repr\_seed0/rotsens\_temporal\_full.png}}
\end{figure}
\begin{figure}[H]
\centering
\includegraphics[width=0.75\textwidth]{plots_small_t30/repr_seed0/rotsens_temporal_full.png}
\caption{Regime B, \FullSpiking, seed 0 --- same quantity, $T{=}30$.
\texttt{plots\_small\_t30/repr\_seed0/rotsens\_temporal\_full.png}}
\end{figure}

\FloatBarrier
\textit{Layer 7 (full capacity) plateaus visibly higher than other layers in \Alternating\ and
\FullSpiking, consistent with Section~\ref{sec:part3}; in \HybridLate\ the plateau levels across layers
are closer together, consistent with that configuration's flatter layer-wise sensitivity profile.
\Analog\ has no meaningful 5.4 plot at full capacity ($T{=}1$, no timestep axis). The plateau-not-decay
result reproduces exactly at reduced capacity, in both regimes, with no exceptions found in any
configuration or seed examined.}

\FloatBarrier

% ===========================================================================
\section{Direct Equivariance Error (Experiment 5.8 / Experiment 2's $E_{90}$)}
\label{sec:part5}

This is the section of the project where the most care is required, because two genuinely different
quantities were measured under the same name (``equivariance error'') across reports, and one of them
(the reduced-capacity $E_{90}$ value) needed a genuine correction. It is deliberately structured in four
steps so none of that gets blurred together.

\subsection{Two Different Quantities, Not Directly Comparable}

Both quantities are legitimate, correctly-implemented computations of the relative error
$\|f(R_\theta x) - f(x)\| / \|f(x)\|$, which the $p4$ architecture guarantees is exactly 0 at the four
group angles $\{0^\circ,90^\circ,180^\circ,270^\circ\}$ by construction. They differ in two ways that
turn out to matter a great deal:

\begin{table}[H]
\centering
\begin{tabular}{@{}lp{4.6cm}p{4.6cm}@{}}
\toprule
 & \textbf{Report 1 (full capacity)} & \textbf{Reports 2--3 (reduced capacity, $E_{90}$)} \\
\midrule
Space & raw accumulated-spike-count logits (no softmax) & softmax class-probability vector \\
Angle reduction & peak (max) over all 24 test angles & value at exactly $\theta{=}90^\circ$ \\
\bottomrule
\end{tabular}
\caption{The two ``equivariance error'' formulas compared across reports.}
\end{table}

\textit{The audit confirmed directly (not just by formula inspection) that this difference, not model
capacity, drives the apparently contradictory orderings below: recomputing BOTH formulas on the SAME
reduced-capacity checkpoints reproduces Report~1's ordering under the raw-logit/peak-over-angles formula
(\None\ highest at 0.417, \FullSpiking\ lowest at 0.270) and reproduces the opposite ordering under the
softmax/value-at-$90^\circ$ formula (\FullSpiking\ highest at 0.498) --- on identical models. Softmax
compresses outputs onto a bounded probability simplex, which changes both which angle maximizes the
disagreement and how the denominator scales with each configuration's confidence/calibration, which
differs systematically by configuration. Sections~\ref{sec:5.2}--\ref{sec:5.3} below report each
quantity as its own report found it; they are not reconciled into one ordering (Section~\ref{sec:5.4}).}

\subsection{Full-Capacity Finding: Peak Equivariance Error (Report 1, As-Is)}
\label{sec:5.2}

Quantity: raw-logit relative error, peak over all 24 angles. Every configuration's error touches (or
comes extremely close to) 0 at exactly the four group angles and rises to a local maximum roughly
halfway between them, in every configuration and every seed --- the expected sawtooth, and a strong
sanity check that \texttt{P4BatchNorm2d}'s equivariance fix functions correctly in fully trained models.

\begin{figure}[H]
\centering
\includegraphics[width=0.75\textwidth]{plots/analog_seed0/analog_seed0_5.8_equivariance_error.png}
\caption{Full capacity, \Analog, seed 0 --- equivariance error vs.\ angle.
\texttt{plots/analog\_seed0/analog\_seed0\_5.8\_equivariance\_error.png}}
\end{figure}
\begin{figure}[H]
\centering
\includegraphics[width=0.75\textwidth]{plots/full_seed0/full_seed0_5.8_equivariance_error.png}
\caption{Full capacity, \FullSpiking, seed 0 --- same quantity.
\texttt{plots/full\_seed0/full\_seed0\_5.8\_equivariance\_error.png}}
\end{figure}

\FloatBarrier
\begin{table}[H]
\centering
\begin{tabular}{@{}llc@{}}
\toprule
\textbf{Configuration} & \textbf{Spiking layers} & \textbf{Approx.\ peak equiv.\ error (mean over seeds)} \\
\midrule
\Analog & 0/7 & $\sim 0.58$--0.59 (highest) \\
\Alternating & 3/7 & $\sim 0.50$--0.60 \\
\HybridLate & 1/7 & $\sim 0.35$--0.40 \\
\FullSpiking & 7/7 & $\sim 0.33$--0.42 (lowest) \\
\bottomrule
\end{tabular}
\caption{Peak equivariance error, full capacity.}
\end{table}

If spiking dynamics simply degraded invariance in proportion to how many layers spike, peak error should
increase monotonically from \Analog\ through \HybridLate\ through \Alternating\ through \FullSpiking.
The data does not show this: \Analog\ has the highest peak error and \FullSpiking\ the lowest --- the
opposite of the naive ``spiking hurts invariance'' expectation. Report 1 reports this as an open,
reproducible observation, not an explained result; no experiment run in that report isolates why this
ordering holds.

\subsection{Reduced-Capacity $E_{90}$ Finding, CORRECTED (Supersedes Reports 2 and 3)}
\label{sec:5.3}

\textit{Quantity: softmax-probability relative error, value at exactly $\theta{=}90^\circ$ (``$E_{90}$'').}

This supersedes two earlier readings of the same quantity. Report 2's single-seed-0 reading first
suggested \FullSpiking\ was uniquely, severely anomalous at the group angles; Report 2's seed-1 revision
(\texttt{findings\_report\_t30.tex}) walked that back to ``noisy across every configuration, no
configuration reproducibly anomalous, $n{=}2$ is not enough to tell.'' Report 3, using a probe set held
fixed across seeds (a genuine methodological improvement over Report 2's per-seed-varying probe set),
found a clean monotonic ordering ($\None < \Alternating < \FullSpiking$) in both regimes and described
it as materially cleaner than Report 2's reading. The correction described in Section~\ref{sec:corrections}
(Fix 2: 8 independent encoding draws per model seed, each with a genuinely different encoding seed) was
applied on top of Report 3's fixed-probe-set protocol, because the audit traced the underlying noise
source to Monte-Carlo spike-sampling variance in the encoding, not to the probe set.

\begin{table}[H]
\centering
\begin{tabular}{@{}lccc@{}}
\toprule
\textbf{Regime} & \None & \Alternating & \FullSpiking \\
\midrule
A ($T{=}50$) & $0.0038 \pm 0.0044$ & $0.0102 \pm 0.0038$ & $0.0249 \pm 0.0151$ \\
B ($T{=}30$) & $0.0146 \pm 0.0087$ & $0.0248 \pm 0.0147$ & $0.0455 \pm 0.0168$ \\
\bottomrule
\end{tabular}
\caption{$E_{90}$, mean $\pm$ std over 3 model seeds; each seed's own value is itself the mean of 8
independent encoding draws, so this combines model-seed variance and encoding-draw variance.}
\label{tab:e90corrected}
\end{table}

\begin{figure}[H]
\centering
\includegraphics[width=0.85\textwidth]{plots_small/experiment2/exp2_e90_multidraw_comparison.png}
\caption{$E_{90}$ before (single-draw) vs.\ after (8-draw mean) the Fix 2 correction, both regimes, mean
$\pm$ std over 3 model seeds. \texttt{plots\_small/experiment2/exp2\_e90\_multidraw\_comparison.png}}
\end{figure}

\FloatBarrier
\textbf{What survives:} the mean ordering $\None < \Alternating < \FullSpiking$ holds in both regimes,
independently --- and it is now corroborated two ways rather than one, since it reproduces at $T{=}50$
and $T{=}30$ on entirely different checkpoints and is stable to which encoding draw is used, not just to
which model seed happened to be picked. That is a real strengthening of the claim relative to any
single-draw reading.

\textbf{What does not survive:} the word ``clean.'' The mean $\pm$ std bands overlap substantially
between adjacent configurations at $n{=}3$ model seeds in both regimes (e.g.\ Regime A: \None's upper
band 0.0083 already overlaps \Alternating's lower band 0.0064; \Alternating/\FullSpiking\ overlap even
more heavily). This is not new noise introduced by the correction --- the old single-draw per-seed
numbers already had the same overlap (e.g.\ Regime A seed 1: $\Alternating{=}0.0213 > \FullSpiking{=}0.0177$)
--- multi-draw averaging simply makes the overlap explicit and quantified rather than hidden behind a
single point estimate per seed. \emph{Directional and cross-regime-consistent} is a defensible claim at
$n{=}3$ seeds; \emph{clean statistical separation} between adjacent configurations is not, and would
require more model seeds, not more encoding draws.

\subsection{These Two Findings Are Not Reconciled}
\label{sec:5.4}

\begin{keypoint}
Section~\ref{sec:5.2} (full-capacity, peak-over-angles, raw logits) and Section~\ref{sec:5.3}
(reduced-capacity, value-at-$90^\circ$, softmax, multi-draw-corrected) are different measurements on
different models. No attempt is made here to merge them into one ordering or one causal story, because
Section~\ref{sec:part5}.1 already demonstrated that switching the formula alone flips the ordering on
identical checkpoints --- any apparent agreement or disagreement between Sections~\ref{sec:5.2}
and~\ref{sec:5.3} as currently computed would not be informative. This is flagged as an explicit open
question for future work: computing BOTH quantities (raw-logit peak-over-angles, and softmax multi-draw
value-at-$90^\circ$) on the SAME set of checkpoints, ideally the full-capacity model, would be needed
before any claim relating them could be made.
\end{keypoint}

\FloatBarrier

% ===========================================================================
\section{Layer/Time Representation Distance $D_{\ell,t}$ (Experiment 3)}
\label{sec:part6}

$D_{\ell,t}(\theta) = \mathbb{E}_x\!\left[\dfrac{\|h_{\ell,t}(R_\theta x) - h_{\ell,t}(x)\|}
{\|h_{\ell,t}(x)\| + \epsilon}\right]$, the representation distance at a single layer and single
timestep, evaluated at $\theta{=}45^\circ$ (no exact group guarantee) and $\theta{=}90^\circ$ (exact
$C_4$ guarantee). This section reports the CORRECTED numbers (the correction described in
Section~\ref{sec:corrections}, Fix 1) and states plainly that the original raw values ($\sim 10^5$ to
$\sim 3.4\times10^7$) were floor-division artifacts carrying no signal about representation distance,
not a real result to be interpreted.

\subsection{Layer 5, Mean $D_{\ell,t}(45^\circ)$: Before vs.\ After}

\begin{figure}[H]
\centering
\includegraphics[width=0.85\textwidth]{plots_small/experiment3_fixed/exp3_layer5_before_after_all_combos.png}
\caption{Layer 5 mean $D_{45^\circ}$, all 3 configs $\times$ 3 seeds $\times$ 2 regimes, before (fixed
$\epsilon{=}10^{-8}$, log scale) vs.\ after (layer-relative $\epsilon$ + floor masking) the Fix 1
correction. \texttt{plots\_small/experiment3\_fixed/exp3\_layer5\_before\_after\_all\_combos.png}}
\end{figure}

\FloatBarrier
\begin{table}[H]
\centering
\small
\begin{tabular}{@{}llrrrr@{}}
\toprule
\textbf{Regime} & \textbf{Config} & \textbf{Seed} & \textbf{Mean BEFORE} & \textbf{Mean AFTER} & \textbf{Cells excl.} \\
\midrule
T50 & \None & 0 & $5.2\times10^5$ & $1.03$ & 1.0\% \\
T50 & \None & 1 & $1.3\times10^6$ & $1.13$ & 6.6\% \\
T50 & \Alternating & 0 & $2.1\times10^6$ & $1.16$ & 7.6\% \\
T50 & \FullSpiking & 0 & $8.9\times10^6$ & $1.18$ & 9.5\% \\
T30 & \FullSpiking & 0 & $9.6\times10^6$ & $1.12$ & 12.8\% \\
\bottomrule
\end{tabular}
\caption{Representative rows; full 18-combination table is in \texttt{exp3\_eps\_diagnostics.csv} under
each regime's \texttt{experiment3\_fixed/} folder.}
\end{table}

\textit{Every blown-up value collapses to the 0.85--1.4 range once the floor artifact is removed; seeds
that were already sane (e.g.\ \None\ seed 2: 0.87 before and after) barely move --- confirming the fix
targets exactly the broken cells and does not disturb the ones that were already correct.}

\subsection{Visual Before/After: Heatmaps}

\begin{figure}[H]
\centering
\includegraphics[width=\textwidth]{plots_small/experiment3/exp3_heatmap_full_seed0.png}
\caption{BEFORE (original, \texttt{plots\_small/experiment3/}): \FullSpiking, Regime A, seed 0 ---
$D_{\ell,t}(45^\circ)$ heatmap. The deepest layer's color range is dominated by floor-division spikes.}
\end{figure}
\begin{figure}[H]
\centering
\includegraphics[width=\textwidth]{plots_small/experiment3_fixed/exp3_heatmap_full_seed0.png}
\caption{AFTER (corrected, \texttt{plots\_small/experiment3\_fixed/}): same model/seed/quantity, with
the Fix 1 layer-relative epsilon and floor masking applied.}
\end{figure}

\FloatBarrier
\subsection{Collapse-Frequency-by-Configuration Finding, Corrected Caveat}

Report 3 originally described the layer-5 floor-collapse frequency as ``rare in \None, frequent in
\Alternating, near-constant in \FullSpiking,'' read from single-seed data. With the corrected metric
computed across all 3 seeds and both regimes, the directional claim survives on average but the
``near-constant'' language does not --- it overstates what a single seed showed.

\begin{table}[H]
\centering
\begin{tabular}{@{}lccc@{}}
\toprule
 & \None & \Alternating & \FullSpiking \\
\midrule
T50, mean \% cells excluded ($\theta{=}45^\circ$) & $2.5\% \pm 3.6$ & $3.1\% \pm 4.0$ & $3.8\% \pm 5.0$ \\
T30, mean \% cells excluded ($\theta{=}45^\circ$) & $2.1\% \pm 3.1$ & $0.7\% \pm 1.1$ & $4.7\% \pm 7.0$ \\
\bottomrule
\end{tabular}
\caption{Layer-5 floor-collapse frequency, mean $\pm$ std over 3 seeds.}
\end{table}

\textit{On average, across seeds, \FullSpiking\ does have the highest floor-collapse rate in both
regimes, consistent with more spiking layers upstream propagating more sparsity into the narrow final
layer. But this is not a clean per-seed pattern: \FullSpiking\ seed 0 hits 9.5--12.8\% exclusion while
\FullSpiking\ seed 1 hits only 0.5--0.6\%, comparable to \None. The correct framing is ``higher on
average, consistent with the number of spiking layers, but with seed-to-seed variance as large as the
effect itself'' --- not ``near-constant.'' This is the same seed-variance caution that recurs elsewhere
in this project (Sections~\ref{sec:part9}--\ref{sec:part10}).}

\begin{figure}[H]
\centering
\includegraphics[width=0.65\textwidth]{plots_small/experiment3_fixed/exp3_lines_full_seed0.png}
\caption{Regime A, \FullSpiking, seed 0 (corrected) --- $D_{\ell,t}(45^\circ)$ solid vs.\
$D_{\ell,t}(90^\circ)$ dashed (layer 5 only), log scale.}
\end{figure}

\FloatBarrier
The upstream-sparsity-propagation hypothesis first proposed in Report 3 --- that sparsity introduced by
an upstream spiking layer propagates forward and makes a downstream analog layer's pre-activation more
likely to be exactly zero at a given timestep, rather than the layer's own activation type being what
matters --- remains a plausible, mechanistically motivated reading of the corrected data, but is still
not directly verified (Section~\ref{sec:part11}).

% ===========================================================================
\section{Fourier Analysis of the Rotational Response (Experiment 4)}
\label{sec:part7}

This experiment needs no correction and is the strongest, most seed-stable result in the project. It is
presented here as reported in Report 3, unchanged, as the headline confirmatory result of the whole
programme.

Setup: for \None, \Alternating, \FullSpiking, all 3 seeds, both regimes, the classification margin
$M(\theta) = z_y(R_\theta x) - \max_{c\neq y} z_c(R_\theta x)$ and accuracy $A(\theta)$ are evaluated at
$1^\circ$ resolution (360 angles) on a fixed 32-image probe set, using raw accumulated-spike-count
logits (not softmax). $M(\theta)$ is Fourier-decomposed up to $k_{\max}{=}20$, giving $E_{\text{noninv}}$
(non-constant energy fraction) and $E_{\text{forbidden}}$ (energy at $k$ not divisible by 4 ---
architecturally, this should be $\approx 0$ for an exact $p4$ network).

\begin{figure}[H]
\centering
\includegraphics[width=0.85\textwidth]{plots_small/experiment4/exp4_rotation_response_full.png}
\caption{Regime A, \FullSpiking\ --- classification margin $M(\theta)$, all 3 seeds + mean, $1^\circ$
resolution.}
\end{figure}
\begin{figure}[H]
\centering
\includegraphics[width=0.85\textwidth]{plots_small/experiment4/exp4_fourier_spectrum_full.png}
\caption{Regime A, \FullSpiking\ --- Fourier spectrum magnitude vs.\ $k$; red = allowed ($k$ a multiple
of 4), gray = forbidden.}
\end{figure}

\FloatBarrier
The spectrum is unambiguous: every forbidden frequency sits at a uniform noise floor
($\sqrt{a_k^2+b_k^2} \lesssim 0.03$) indistinguishable from zero, while the allowed frequencies show a
clear, decaying spectrum dominated by $k{=}4$ ($\sim 2$--3), with smaller but clearly nonzero
contributions at $k{=}8,12,16,20$. This is the clearest confirmation in the entire project of the $p4$
architecture's exact discrete-symmetry guarantee, directly in Fourier space rather than through an
accuracy or representation-similarity proxy.

\begin{table}[H]
\centering
\scriptsize
\begin{tabular}{@{}llrrrrr@{}}
\toprule
\textbf{Regime} & \textbf{Config} & $a_0$ & $E_{\text{noninv}}$ & $E_{\text{forbidden}}$ & $|coef_{k=4}|$ & $|coef_{k=8}|$ \\
\midrule
A & \None & $7.11\pm2.62$ & $0.278\pm0.094$ & $9.2\times10^{-5}\pm3.9\times10^{-5}$ & $4.07\pm0.74$ & $0.63\pm0.28$ \\
A & \Alternating & $7.80\pm1.88$ & $0.228\pm0.058$ & $1.11\times10^{-4}\pm5.5\times10^{-5}$ & $4.08\pm0.34$ & $0.36\pm0.34$ \\
A & \FullSpiking & $5.94\pm1.61$ & $0.192\pm0.086$ & $1.18\times10^{-4}\pm3.8\times10^{-5}$ & $2.70\pm0.51$ & $0.48\pm0.11$ \\
B & \None & $5.73\pm0.96$ & $0.252\pm0.095$ & $1.56\times10^{-4}\pm2.6\times10^{-5}$ & $3.31\pm1.16$ & $0.37\pm0.15$ \\
B & \Alternating & $6.27\pm0.99$ & $0.174\pm0.064$ & $1.33\times10^{-4}\pm4.9\times10^{-5}$ & $2.77\pm0.23$ & $0.22\pm0.14$ \\
B & \FullSpiking & $4.07\pm0.66$ & $0.201\pm0.050$ & $2.84\times10^{-4}\pm6.7\times10^{-5}$ & $2.01\pm0.44$ & $0.22\pm0.18$ \\
\bottomrule
\end{tabular}
\caption{Experiment 4 summary, mean $\pm$ std over 3 seeds.}
\end{table}

\begin{figure}[H]
\centering
\includegraphics[width=\textwidth]{plots_small/experiment4/exp4_enoninv_eforbidden.png}
\caption{Regime A --- $E_{\text{noninv}}$ (left) vs.\ $E_{\text{forbidden}}$ (right), boxplot + individual
seeds, all 3 configurations.}
\end{figure}

\FloatBarrier
\begin{keypoint}
$E_{\text{forbidden}}$ is architectural, not learned: all three configurations cluster tightly in the
same narrow band ($\sim 6\times10^{-5}$ to $\sim 2.8\times10^{-4}$) across both regimes, with heavily
overlapping seed distributions --- exactly the ``essentially independent of the random seed''
prediction this experiment was designed to test. $E_{\text{noninv}}$ is NOT architectural: it varies
more across configurations and seeds (0.17--0.28), with a mild trend toward lower non-invariant energy
as more layers spike, visible for the $k{=}4$ coefficient specifically ($\None > \Alternating >
\FullSpiking$ in both regimes) though not fully consistent for $E_{\text{noninv}}$ itself. This is
exactly the architecture-determines-symmetry / training-determines-learned-invariance-strength
distinction the experimental programme set out to isolate.
\end{keypoint}

% ===========================================================================
\section{Compute-Budget Findings: Truncated-$T'$ Accuracy (Experiment 5.7)}
\label{sec:part8}

A single model, trained once at its regime's full $T$, has its recorded spike train truncated to the
first $T'$ timesteps at evaluation time and a classification decision decoded from only those, without
retraining. This tests whether rotation-robust accuracy structure emerges progressively as spikes
accumulate, and how quickly.

\subsection{Full Capacity (Report 1)}

\begin{figure}[H]
\centering
\includegraphics[width=0.75\textwidth]{plots/full_seed0/full_seed0_5.7_temporal_evolution.png}
\caption{Full capacity, \FullSpiking, seed 0 --- accuracy vs.\ angle at $T'{=}1,5,10,20,50$.}
\end{figure}
\begin{figure}[H]
\centering
\includegraphics[width=0.75\textwidth]{plots/hybrid_late_seed0/hybrid_late_seed0_5.7_temporal_evolution.png}
\caption{Full capacity, \HybridLate, seed 0 --- same $T'$ grid.}
\end{figure}

\FloatBarrier
\FullSpiking's $T'{=}1$ curve is flat at $\sim 25$--27\% (chance) across every angle --- with all 7
layers spiking, every layer needs several timesteps to integrate enough input spikes to pass a
meaningful signal onward, and this startup delay compounds across all 7 layers. Real W-shaped structure
appears only by $T'{=}10$ and is nearly indistinguishable from $T'{=}50$ by $T'{=}20$. \Alternating's
$T'{=}1$ curve already shows a weak but genuine W-shape (25--70\%, seed/angle-dependent). \HybridLate's
$T'{=}1$ curve already closely tracks the full $T'{=}50$ curve's shape (63--99\%) --- with only the final
layer spiking, six of seven layers process instantaneously regardless of $T'$. This gives a consistent,
monotonic ordering across all 9 seed-runs examined (3 configs $\times$ 3 seeds): the more layers that
spike, the more timesteps needed before rotation-robust accuracy structure appears.

\subsection{Reduced Capacity, Both Regimes (Report 2)}

\begin{table}[H]
\centering
\small
\begin{tabular}{@{}llrrrrr@{}}
\toprule
\textbf{Regime} & \textbf{Config} & $T'{=}1$ & $T'{=}5$ & $T'{=}10$ & $T'{=}20$ & $T'{=}50$ \\
\midrule
A & \Alternating & 51.2\% & 84.4\% & 86.6\% & 87.0\% & 89.4\% \\
A & \FullSpiking & 31.2\% & 71.0\% & 82.6\% & 84.7\% & 85.9\% \\
A & \HybridLate & 58.7\% & 87.0\% & 88.7\% & 89.1\% & 89.6\% \\
A & \None & 77.1\% & 85.9\% & 88.4\% & 88.7\% & 90.5\% \\
\midrule
B & \Alternating & $45.5\pm6.6\%$ & $79.6\pm12.6\%$ & $83.9\pm9.7\%$ & $85.1\pm9.0\%$ & $84.9\pm9.0\%$ \\
B & \FullSpiking & $34.3\pm6.9\%$ & $66.1\pm17.4\%$ & $78.6\pm8.9\%$ & $82.0\pm4.8\%$ & $83.3\pm4.3\%$ \\
B & \HybridLate & $60.3\pm9.1\%$ & $76.0\pm7.3\%$ & $79.1\pm6.0\%$ & $80.3\pm6.7\%$ & $81.5\pm6.3\%$ \\
B & \None & $73.1\pm6.1\%$ & $80.9\pm4.0\%$ & $83.2\pm3.7\%$ & $84.5\pm3.1\%$ & $84.8\pm3.0\%$ \\
\bottomrule
\end{tabular}
\caption{Truncated-$T'$ accuracy, reduced capacity. Regime A ($T{=}50$): single seed (seed 0). Regime B
($T{=}30$): mean $\pm$ std over seeds $\{0,1\}$.}
\end{table}

\begin{figure}[H]
\centering
\includegraphics[width=0.75\textwidth]{plots_small/repr_seed0/truncated_T_by_angle_full.png}
\caption{Regime A, \FullSpiking, seed 0 --- accuracy vs.\ angle at $T'{=}1,5,10,20,50$.}
\end{figure}
\begin{figure}[H]
\centering
\includegraphics[width=0.75\textwidth]{plots_small/repr_seed0/truncated_T_by_angle_hybrid_late.png}
\caption{Regime A, \HybridLate, seed 0 --- same $T'$ grid.}
\end{figure}

\FloatBarrier
The same ``more spiking layers $\Rightarrow$ longer warm-up'' ordering reproduces at reduced capacity,
in both regimes.

\begin{keypoint}
Practical implication, consistent across every capacity and regime tested: $T'{=}10$ already recovers
most of the $T'{=}50$ accuracy everywhere (full capacity: visually indistinguishable by $T'{=}20$;
reduced capacity Regime A: $\geq 94\%$ of $T'{=}50$ accuracy for every configuration; reduced capacity
Regime B: 93--99\% of the $T{=}30$-clamped $T'{=}50$ accuracy for every configuration). Models trained
and evaluated at the full $T$ throughout this project may be using substantially more simulation time
than necessary at decode time.
\end{keypoint}

\FloatBarrier

% ===========================================================================
\section{What Was Ruled Out: Code vs.\ Methodology Audit}
\label{sec:part9}

This is a permanent settled-questions record of every anomaly investigated during this project's
code-vs-methodology audit, so a future reader does not need to re-investigate them. For each: what was
suspected, what was ruled out, the verdict (code bug vs.\ methodological/metric artifact), and the
one-line resolution.

\subsection{\texttt{P4BatchNorm2d} Breaking Equivariance After Training (Resolved Before This Audit)}

\textit{Suspected: standard BatchNorm's independent per-orientation-channel statistics might interact
badly with upright-only training.}

\textbf{Verdict: CONFIRMED CODE ISSUE, already fixed.} Untrained equivariance error $\sim 10^{-6}$;
after training with standard BatchNorm, $\sim 0.1$--0.2. Resolution: \texttt{P4BatchNorm2d}, sharing
statistics/affine parameters across the 4 orientation channels, restores $\sim 10^{-6}$ post-training.
Used throughout every result in this document. This is the one genuine architecture-level bug found and
fixed in this project.

\subsection{Group-Angle Equivariance-Error Noise (Reports 2--3)}

\textit{Suspected: at exact group angles ($90^\circ/180^\circ/270^\circ$), equivariance error should be
$\approx 0$ by construction, but was observed 10--100$\times$ higher than expected for some configs,
inconsistently across seeds.}

\begin{itemize}[leftmargin=1.6em]
\item \textbf{Ruled out (1):} Rotation-implementation noise (\texttt{grid\_sample} vs.\
\texttt{torch.rot90}). Max abs pixel difference $\sim 2\times10^{-6}$ (float32 noise floor), 4--5 orders
of magnitude below the observed 0.01--0.06 errors; feeding the model the bit-exact rotation instead left
the error unchanged to 6 decimal places.
\item \textbf{Ruled out (2):} Mode-dependent code branching in \texttt{gconv\_small.py}. \texttt{P4ConvZ2},
\texttt{P4ConvP4}, \texttt{P4BatchNorm2d}, \texttt{P4GroupPool} contain no reference to \texttt{mode} or
\texttt{spiking\_layers} anywhere --- one code path for every configuration, confirmed by direct code
reading, not inference.
\end{itemize}

\textbf{Verdict: METHODOLOGICAL ARTIFACT, not a code bug.} The evaluation code resets the encoding RNG
to the same seed value for $x$ and $R_\theta x$. Because rotation moves image content between tensor
positions, this reproduces the same raw random stream but not a matched spike encoding under rotation
--- it gives run-to-run reproducibility, not the noise-cancellation the convention seems to have
intended. Confirmed empirically: re-drawing the encoding seed alone (same model, same probe images, no
rotation or config change) reproduces the exact reported magnitudes and $\None < \Alternating <
\FullSpiking$ ordering. Also discovered along the way: layer 1's input is stochastically spike-encoded
in EVERY configuration, including \None --- ``Analog'' is not literally analog at the input, only in
its internal layers --- which is why \None\ also showed nonzero, seed-inconsistent noise at group angles
rather than being immune to the effect. Resolution: Fix 2 (Section~\ref{sec:corrections}), multi-draw
averaging.

\subsection{$D_{\ell,t}$ Reference-Norm Collapse (Report 3, Experiment 3)}

\textit{Suspected: the deepest layer's $D_{\ell,t}(45^\circ)$ values explode to $\sim 10^7$ because the
denominator $\|h_{\ell,t}(x)\| + \epsilon$ is near-zero when that timestep's activation is mostly/all
zero.}

\textbf{Verdict: METHODOLOGICAL/METRIC-DESIGN ARTIFACT, not a model or training bug.} $\epsilon{=}10^{-8}$
was a fixed global constant; measured layer-5 reference norms directly on real checkpoints and confirmed
9.5\% of (timestep, image) cells there have a raw norm of EXACTLY zero (the whole 16-unit layer silent
that timestep), so $\text{denom}{=}\epsilon$ exactly and any nonzero numerator explodes the ratio.
Layers 1--4 never hit an exact zero in the same probe set --- the problem is specific to the narrowest,
sparsest layer. Resolution: Fix 1 (Section~\ref{sec:corrections}), layer-relative epsilon plus floor
masking.

\subsection{Report 1 vs.\ Report 3 Equivariance-Error Direction Reversal}

\textit{Suspected: Report 1 found peak equivariance error highest for \Analog, lowest for \FullSpiking;
Reports 2--3 found error at $90^\circ$ specifically lowest for \Analog/\None, highest for \FullSpiking
--- apparently contradictory.}

\textbf{Verdict: NOT A BUG IN EITHER} --- genuinely different quantities being compared (see
Section~\ref{sec:part5}.1 for the full formula comparison). Confirmed directly: recomputing both
formulas (raw-logit peak-over-angles; softmax value-at-$90^\circ$) on the IDENTICAL reduced-capacity
checkpoints reproduces each report's ordering under its own formula --- the formula choice alone flips
the ranking, not model capacity. Both implementations correctly compute what they claim to; the issue is
that the two numbers were juxtaposed across reports as if directly comparable ``equivariance error''
when they are not. Left unreconciled by design (Section~\ref{sec:5.4}).

\subsection{Evolution of the ``\FullSpiking\ Anomaly'' Narrative Across Reports}

For completeness: the understanding of Section~9.2 evolved visibly across this project's own reports,
and this evolution is itself worth recording so it is not mistaken for inconsistency.
\texttt{findings\_report\_small.tex} (seed-0-only reading) proposed ``cascaded LIF-threshold
amplification of spike-sampling stochasticity, scaling with spiking-layer count'' as the explanation,
based on \FullSpiking\ looking uniquely and reproducibly elevated across its two available regimes.
\texttt{findings\_report\_t30.tex} (after seed 1 arrived) walked this back: four of five configurations
showed at least one elevated group-angle value depending on seed, so no single configuration could
responsibly be called uniquely anomalous at $n{=}2$. The final audit (Section~9.2, this document)
confirms the original mechanism's core intuition was directionally right (more spiking layers does
correlate with more noise, confirmed cleanly via the multi-draw experiment in Section~\ref{sec:5.3}) but
locates the root cause more precisely (RNG/pixel-position mismatch at the stochastically-encoded input,
present in every configuration including \None, not a LIF-threshold-specific cascade) and treats the
effect as quantifiable Monte-Carlo variance rather than a per-configuration ``bug.'' Later readings are
more reliable than earlier ones here; this document uses the final audit's framing throughout.

% ===========================================================================
\section{Synthesis}
\label{sec:part10}

Written at the confidence level each finding actually supports --- Section~\ref{sec:5.3}'s now-corrected
$E_{90}$ finding is not upgraded back to ``clean ordering'' language here just because this is the
summary.

\begin{enumerate}[leftmargin=1.8em]
\item \textbf{Accuracy alone is uninformative about spike placement}, at every capacity and regime
tested. Once training is done properly, all configurations reach rotation-accuracy curves within a few
points of each other, with one reproducible exception (\HybridEarly's large seed-to-seed variance at
reduced capacity, absent at full capacity only because Report 1 never trained it).
\item \textbf{The layer-wise representational collapse at the deepest layer(s) is a genuine, unhedged,
universal finding}: it appears in every configuration at every capacity tested, including the
zero-spiking control, and is a property of the $p4$-equivariant architecture's final stages, not of
spiking dynamics.
\item \textbf{Rotation sensitivity through spiking time rises quickly and then plateaus} --- it does not
decay --- in every configuration, every seed, and every regime tested, directly contradicting the
project's original temporal-contraction hypothesis. Per-configuration differences in the shape of the
$\mathcal{R}_i$-vs-layer climb (sharp late jump, gradual climb, or flat) are real at full capacity and
reproduce qualitatively but less cleanly at reduced capacity, where 5 layers give less room to resolve
the shapes distinctly.
\item \textbf{Truncated-$T'$ accuracy shows a clean, reproducible, mechanistically sensible ordering}
across every capacity and regime: more spiking layers means a longer warm-up before rotation-robust
accuracy structure appears, and $T'{=}10$ recovers most of the full-$T$ accuracy everywhere --- a
genuine, actionable compute-budget finding.
\item \textbf{Direct equivariance error tells two different, non-reconciled stories} depending on
exactly what is measured. At full capacity, peak error over all angles is non-monotonic in spiking
extent (\Analog\ worst, \FullSpiking\ best) and remains unexplained. At reduced capacity, the corrected
(multi-draw) $E_{90}$ at exactly $90^\circ$ shows a directionally consistent, cross-regime-reproducible
ordering ($\None < \Alternating < \FullSpiking$ in the mean) that does NOT achieve clean separation
between adjacent configurations at $n{=}3$ model seeds --- a real but modest finding, not the ``clean
ordering'' one of its source reports originally claimed.
\item \textbf{Experiment 4's Fourier analysis is the project's strongest and most seed-stable
confirmatory result}: forbidden-mode energy is architecturally pinned near zero regardless of
configuration or seed, while allowed-mode amplitude (the learned component) varies substantially --- a
clean empirical separation of what the architecture guarantees from what training actually learns.
\item \textbf{Two of the three anomalies investigated in this project's code-vs-methodology audit were
genuine methodological artifacts} (Monte-Carlo spike-encoding noise; a fixed epsilon inappropriate for a
narrow layer), now corrected in the analysis code; none was a bug in the $p4$-equivariant conv/BatchNorm
implementation itself, which was independently verified to have no mode-dependent branching. The one
genuine historical code bug in this project (standard BatchNorm breaking equivariance) was found and
fixed before this audit began.
\item \textbf{A recurring methodological lesson across this entire project}, worth stating on its own:
single-seed or single-draw readings repeatedly looked cleaner than they turned out to be once a second
seed, a second draw, or a corrected metric arrived (the \FullSpiking-anomaly narrative, the
$D_{\ell,t}$ collapse-frequency ordering, and the $E_{90}$ ordering all followed this pattern). Every
finding in this document that still holds at $n{=}3$ seeds has been cross-checked against this pattern
before being stated plainly.
\end{enumerate}

% ===========================================================================
\section{Limitations and Recommended Next Steps}
\label{sec:part11}

\subsection{Limitations, Consolidated Across All Reports}

\begin{itemize}[leftmargin=1.8em]
\item \textbf{Seed counts below programme recommendations.} Every seed-averaged number in this document
uses $n{=}3$ model seeds ($n{=}2$ for several reduced-capacity Regime B representation-level quantities,
since seed 2's representation analysis was intentionally left incomplete by explicit instruction). The
Priority Programme recommends $n{=}10$ (Experiment 2) and $n{=}5$ (Experiment 3); this project has not
reached either.
\item \textbf{SVCCA absolute-scale caveat.} SVCCA is not on an absolute 0--1 similarity scale ---
independent, genuinely dissimilar data produces a nonzero baseline of roughly 0.2--0.4 at the batch
sizes used here. Read SVCCA values relatively (larger = more similar), never as an absolute invariance
measure.
\item \textbf{$\mathcal{R}_i$ cross-activation-type scale caveat.} $\mathcal{R}_i$'s absolute magnitude
depends on whether the underlying activation is a continuous ReLU output or a spike count (the spiking
configurations' $\mathcal{R}_i$ sits roughly $1000\times$ larger than \Analog/\None's at comparable
layers). Only within-configuration comparisons (across layers or seeds) are on safe footing;
cross-configuration comparisons must use shape, not magnitude.
\item \textbf{Untrained \HybridEarly/\HybridMid\ at full capacity.} Report 1 never trained \HybridEarly\
or \HybridMid, so Sections~\ref{sec:part3}/\ref{sec:part4}/\ref{sec:part5}/\ref{sec:part8}'s
full-capacity subsections say nothing about them; only the reduced-capacity reports (which trained
\HybridEarly, and have no \HybridMid\ analogue at depth 5) cover it.
\item \textbf{Experiment 1 not run.} A parameter/depth-matched non-equivariant SNN baseline was never
run at any capacity --- it requires building and training a new architecture, out of scope for every
pass in this project so far, all of which were either representation-analysis-only or explicitly
no-retraining.
\item \textbf{Probe-set quantization.} Report 3's 32-image fixed probe set quantizes accuracy-derived
quantities ($D_s$ in particular) in steps of $1/32 \approx 3.1\%$; this is why Experiment 2's $D_s$
should be read as indicative, not precise, at this probe-set size.
\item \textbf{$D_{\ell,t}$ 45-vs-90 comparison still limited at layer 5.} The Fix 1 correction addresses
the reference-norm-collapse artifact, but Experiment 3's original design goal (cleanly comparing
$D_{\ell,t}$ at $\theta{=}45^\circ$ vs.\ $\theta{=}90^\circ$ to show the exact-group-angle guarantee
reducing sensitivity specifically at $90^\circ$) is still not achievable at the deepest layer even
post-fix, since both angles are compared against the same reference representation and are equally
exposed to any residual near-zero-norm cells.
\item \textbf{Upstream-sparsity-propagation mechanism unconfirmed.} The claim that sparsity from an
upstream spiking layer propagates into a downstream analog layer's collapse frequency
(Section~\ref{sec:part6}.3) is a plausible, mechanistically motivated reading of the corrected data, not
a directly verified mechanism (e.g.\ via direct inspection of layer 4's pre-layer-5 output sparsity per
configuration).
\item \textbf{\HybridEarly\ seed-2 instability uninvestigated.} \HybridEarly's large seed-2 accuracy
instability (Section~\ref{sec:part2}.2--2.3) was noted in both reduced-capacity regimes but never
investigated further --- unclear if it is a property of that configuration at this capacity or an
isolated bad initialization.
\end{itemize}

\subsection{Recommended Next Steps, Prioritized Against the Sept 19 Abstract / Sept 24 Full-Paper Deadlines}

\textit{Today's date: August 31, 2026 --- roughly 19 days to the abstract deadline, 24 to the full paper.}

\begin{enumerate}[leftmargin=1.8em]
\item \textbf{Immediate (before Sept 19), zero additional compute:} both corrections (Fix 1, Fix 2) are
already applied and this document already reflects them --- no further action needed before the
abstract; simply draw the abstract's claims from Sections~\ref{sec:5.3}, \ref{sec:part6}, and
\ref{sec:part7} of this document rather than the superseded single-draw/fixed-epsilon numbers in the
original Report 2/3 files.
\item \textbf{Cheap, before Sept 24:} directly test the upstream-sparsity-propagation mechanism
(Section~\ref{sec:part6}.3/\ref{sec:part11}.1) by inspecting layer 4's pre-layer-5 output sparsity per
configuration on existing checkpoints --- cheap, no retraining, would upgrade a plausible reading to a
confirmed one before the full paper.
\item \textbf{Moderate effort, before Sept 24 if time allows:} if GPU time allows before the full-paper
deadline, add 2 more model seeds (to $n{=}5$) for Experiment 2's $E_{90}$ specifically, since
Section~\ref{sec:5.3} identified overlapping mean$\pm$std bands between adjacent configurations as the
main remaining weakness in an otherwise cross-regime-consistent finding --- this is the single
highest-leverage additional-compute item in the project given it directly strengthens the report's most
novel corrected finding.
\item \textbf{Moderate effort:} complete reduced-capacity Regime B seed 2's representation-level
analysis (a few-minutes-per-stage job) to bring Sections~\ref{sec:part4}/\ref{sec:part6}/\ref{sec:part8}'s
Regime B tables from $n{=}2$ to $n{=}3$, matching Regime A.
\item \textbf{Larger effort, only if time remains after the above:} train \HybridEarly\ and \HybridMid\
at full capacity, extending Sections~\ref{sec:part3}--\ref{sec:part5}/\ref{sec:part8}'s full-capacity
subsections to the same 5-configuration coverage the reduced-capacity reports already have; run a
controlled ablation isolating which of \texttt{P4BatchNorm2d}/longer training/higher $T$ explains the
original compute-constrained draft's non-replicating findings (Section~\ref{sec:part2}.1).
\item \textbf{Largest effort, full-paper timeframe only:} Experiment 1 (parameter/depth-matched
non-equivariant SNN baseline) --- the largest remaining lift in the project, requiring a new
architecture and full retraining; realistically a full-paper-cycle item rather than an abstract-cycle
one.
\end{enumerate}

\end{document}
